% appendix-synthobj.tex
% Top-level, \input-able fragment documenting the synthobj package. Imposes
% nothing on its host: no \documentclass, no \begin{document}, no \cite, no
% hyperref (the host loads that last, if at all).
%
% Three host hooks, all defined in synthobj-notation.sty and overridable by
% \renewcommand BEFORE this file is \input:
%   \SOlevelA / \SOlevelB / \SOlevelC  -- sectioning; default to
%       \section/\subsection/\subsubsection. A chapter-based thesis (e.g.
%       the Aalto template) renews them to \chapter/\section/\subsection.
%   \SOroot  -- directory holding this fragment's tables/; default ".". A
%       host that relocates docs/thesis/ elsewhere renews \SOroot instead
%       of editing every \input below.
%   \SOciteOAK / \SOciteSAASBO / \SOcitePapenmeier  -- plain-text citation
%       stand-ins; a host with a bibliography renews each to \cite{key}.
\SOlevelA{Synthetic objectives (\SOcode{synthobj})}\label{app:so:top}
% tables/numbers.tex holds one \newcommand{\SOnum<Name>} per measured number
% (generated by regen_tables.py). It is input exactly once, here, so every
% fragment below can use the macros; a host that relocates docs/thesis/
% needs only \SOroot, exactly as for every other table input.
\input{\SOroot/tables/numbers}
\input{\SOroot/A1-construction}
\input{\SOroot/A2-families}
\input{\SOroot/A3-api-fileformat}
\input{\SOroot/A4-design-decisions}
